\documentclass[a4paper]{article}
% Español (México) con XeLaTeX: fontspec para fuentes Unicode, polyglossia
% para la división silábica y los nombres del documento en la variante mexicana.
\usepackage{fontspec}
\usepackage{polyglossia}
\setdefaultlanguage[variant=mexican]{spanish}
% Los nombres chinos van como «pinyin (original)»: xeCJK da a los caracteres
% Han su propia fuente; sin ella XeLaTeX los omite con un aviso.
% FandolSong no cubre algunos caracteres de nombres (U+73FA); WenQuanYi Zen Hei sí.
\usepackage[AutoFallBack=true]{xeCJK}
\setCJKmainfont{FandolSong-Regular.otf}
\setCJKfallbackfamilyfont{\CJKrmdefault}{WenQuanYi Zen Hei}
% polyglossia no traduce los nombres de listings.
\AtBeginDocument{\providecommand{\lstlistingname}{}\renewcommand{\lstlistingname}{Listado}%
  \providecommand{\lstlistlistingname}{}\renewcommand{\lstlistlistingname}{Índice de listados}}
\usepackage{amsmath, amssymb}
\usepackage{graphicx}
\usepackage[margin=2.5cm]{geometry}
\usepackage[most]{tcolorbox}
\usepackage{etoolbox}
\usepackage{listings}
\usepackage{hyperref}
\usepackage{booktabs}
\usepackage{subcaption}
\usepackage{float}
\usepackage{array}

\newtcolorbox{knowledgebox}[1]{enhanced, colback=blue!5!white, colframe=blue!75!black, colbacktitle=blue!75!black, coltitle=white, fonttitle=\bfseries, title={#1}, attach boxed title to top left={yshift=-2mm, xshift=2mm}, boxrule=1pt, sharp corners}
\newtcolorbox{importantbox}[1]{enhanced, colback=yellow!10!white, colframe=yellow!80!black, colbacktitle=yellow!80!black, coltitle=black, fonttitle=\bfseries, title={#1}, sharp corners}
\newtcolorbox{warningbox}[1]{enhanced, colback=red!5!white, colframe=red!75!black, colbacktitle=red!75!black, coltitle=white, fonttitle=\bfseries, title={#1}, sharp corners}

\lstset{
    language=Python,
    basicstyle=\ttfamily\small,
    keywordstyle=\color{blue},
    stringstyle=\color{red!60!black},
    commentstyle=\color{green!60!black},
    breaklines=true,
    frame=single,
    numbers=left,
    numberstyle=\tiny\color{gray},
    captionpos=b,
    extendedchars=true
}

\newcommand{\notetitle}{[LLM Agents SP25] Code Agents \\ \& AI Vulnerability Detection}
\newcommand{\noteauthors}{Elaboradas a partir de la clase de Charles Sutton}
\newcommand{\notedate}{2026-04-02}
\newcommand{\videochannel}{Berkeley RDI}
\newcommand{\videopublishdate}{2025-03-10}
\newcommand{\videoduration}{1:27:02}
\newcommand{\videourl}{https://www.youtube.com/watch?v=JCk6qJtaCSU}
\newcommand{\videocoverpath}{cover.jpg}
\newcommand{\repourl}{https://github.com/hqhq1025/ai-course-notes}

% Footer with repo info
\usepackage{fancyhdr}
\pagestyle{fancy}
\fancyhf{}
\fancyfoot[C]{\thepage}
\fancyfoot[R]{\footnotesize\color{gray}\href{\repourl}{AI Course Notes}}
\renewcommand{\headrulewidth}{0pt}
\renewcommand{\footrulewidth}{0pt}

\begin{document}
\begin{titlepage}
\centering
{\Large Notas del curso\par}
\vspace{1.2cm}
{\huge\bfseries \notetitle\par}
\vspace{0.8cm}
{\large \noteauthors\par}
\vspace{0.3cm}
{\large \notedate\par}
\vspace{1.2cm}
\ifdefempty{\videocoverpath}{}{
\includegraphics[width=0.82\textwidth,height=0.45\textheight,keepaspectratio]{\videocoverpath}\par
}
\vfill
\begin{tcolorbox}[width=0.9\textwidth, colback=black!2!white, colframe=black!60, sharp corners]
\textbf{Autor o canal del video}: \videochannel\par
\textbf{Fecha de publicación}: \videopublishdate\par
\textbf{Duración del video}: \videoduration\par
\ifdefempty{\videourl}{}{\textbf{Enlace del video}: \href{\videourl}{\nolinkurl{\videourl}}\par}
\textbf{Página del proyecto}: \href{\repourl}{\nolinkurl{\repourl}}\par
\end{tcolorbox}
\end{titlepage}
\tableofcontents
\newpage

\section{Ubicación del curso y tesis central}

Esta clase es un tema de alta densidad dentro del curso de primavera «Advanced LLM Agents» de Berkeley, impartido por Charles Sutton, de Google DeepMind. La conferencia parece cubrir dos direcciones (Code Agent y AI Security), pero el hilo conductor real es uno solo: \textbf{cómo llevar al modelo de lenguaje de «saber escribir código» a «poder completar tareas de software abiertas»}. La estructura de la exposición se despliega en tres partes: primero regresa a la evaluación y al diseño de sistemas del coding agent, después discute la adaptación de la IA en tareas de seguridad y, por último, usa el proyecto Big Sleep para explicar por qué el método basado en agentes es válido en la búsqueda real de vulnerabilidades.

\begin{importantbox}{El criterio de ingeniería de toda la clase}
La postura de Sutton no es que «entre más grandes los parámetros del modelo, mejor», sino que «el sistema pueda cerrar el ciclo». Cerrar el ciclo incluye: definición ejecutable de la tarea, resultado verificable, trayectoria iterable, e interfaz de herramientas escalable. Basta con que falte uno de estos cuatro elementos para que el resultado del agente tienda a quedarse en el nivel de demostración.
\end{importantbox}

La primera mitad de la exposición insiste varias veces en un hecho que se pasa por alto con facilidad: desarrollar un agente no es optimizar un prompt una sola vez, sino \textbf{hacer hill climbing en un espacio combinatorio enorme}. Cuando el equipo modifica continuamente la system instruction, la forma de las herramientas, el formato de retroalimentación y la proporción del post-entrenamiento, cada paso depende de la puntuación de la evaluación como guía. Si el objetivo de evaluación está mal alineado, la dirección completa de la iteración se desvía.

\begin{knowledgebox}{Por qué importa el trasfondo interdisciplinario}
El ponente relata primero su propia trayectoria (PLN $\rightarrow$ síntesis de programas $\rightarrow$ generación de código $\rightarrow$ seguridad), no como preámbulo biográfico, sino para dejar claro que: el valor de una tecnología de IA depende de qué tan a fondo se entienda el dominio del problema. Quien solo entiende el modelo, sin conocer la ingeniería de software o los procesos de seguridad, termina optimizando el sistema para que «sepa responder preguntas», no para que «sepa hacer las cosas».
\end{knowledgebox}

\begin{warningbox}{No entiendas esta clase como «más benchmarks de seguridad»}
El punto central de la conferencia no es presentar unas cuantas tablas de clasificación más, sino explicar por qué muchos benchmarks son eficaces en una etapa temprana y dejan de serlo después, y por qué los escenarios de seguridad exigen una verificación dinámica más sólida. Quien solo memorice los nombres de los conjuntos de datos se perderá la metodología de ingeniería más valiosa.
\end{warningbox}

\begin{figure}[H]
\centering
\includegraphics[width=0.88\textwidth]{frame-01-agenda.jpg}
\caption{La agenda en tres partes presentada al inicio de la conferencia: fundamentos del Code Agent, AI Security y el caso Big Sleep \protect\footnotemark}
\end{figure}
\footnotetext{Intervalo de tiempo en el video: 00:03:22--00:03:50.}

\subsection{Resumen de la sección}
Esta sección establece el sistema de referencia para leer toda la clase: no se trata de un repaso disperso de técnicas, sino de una lección de método sobre «cómo evaluar, acotar e implementar correctamente los sistemas de agentes». Todos los detalles posteriores están al servicio de esta pregunta central.
\section{La esencia del Agente: el ciclo ReAct, las llamadas a herramientas y el cómputo en tiempo de prueba}

\subsection{De buzzword a la definición mínima}
Sutton, hacia el minuto 00:04, da directamente su propia definición: lo que llamamos un Agente LLM es, en esencia, \textbf{cómputo dinámico + uso de herramientas}. El cómputo dinámico significa que el modelo no produce la respuesta final de una sola vez, sino que asigna más cómputo en tiempo de prueba según la dificultad de la tarea; el uso de herramientas significa que puede externalizar sus hipótesis intermedias hacia un entorno ejecutable y corregir continuamente su siguiente estrategia mediante comandos, ejecución de código e información de depuración.

\begin{knowledgebox}{El ciclo mínimo de un Agente ejecutable}
\[
\text{Trajectory}_{t}
\xrightarrow{\text{LLM}}
\text{Thought + Tool Call}
\xrightarrow{\text{Environment}}
\text{Observation}_{t+1}
\xrightarrow{\text{append}}
\text{Trajectory}_{t+1}
\]
Este ciclo coincide con la idea de ReAct: el razonamiento no es el punto final, sino lo que permite decidir la siguiente acción con una herramienta.
\end{knowledgebox}

\begin{importantbox}{Por qué el «cómputo dinámico» es más determinante que la «inteligencia de una sola vez»}
\begin{itemize}
    \item Las tareas complejas suelen cometer errores desde los primeros pasos; solo permitir la iteración expone y corrige las rutas equivocadas.
    \item El mismo modelo, bajo ``múltiples rondas de prueba y error + retroalimentación ejecutable'', tiene un techo de capacidad notablemente más alto que bajo ``una sola generación''.
    \item Para problemas de ingeniería reales, la verificación continua vale más que la fluidez de un solo texto.
\end{itemize}
\end{importantbox}

Esta definición coincide con la frase que se repite en la clase: \textit{``Chain of Thought plus tool use is at the heart of agents.''} En escenarios de coding, la salida de la herramienta es la señal de entrenamiento para la siguiente ronda de razonamiento; en escenarios de seguridad, el depurador, el sanitizer y los registros de fallos del programa se convierten en una retroalimentación externa más fuerte.

\begin{figure}[H]
\centering
\includegraphics[width=0.88\textwidth]{frame-07-react-loop.jpg}
\caption{En la parte final del curso se retoma el ciclo de estilo ReAct: observar, analizar, llamar a la herramienta y observar de nuevo \protect\footnotemark}
\end{figure}
\footnotetext{Intervalo de tiempo en el video: 01:18:06--01:18:40.}

\subsection{Presupuesto de cómputo en tiempo de prueba: seguir corrigiendo o reiniciar la trayectoria}
Una pregunta muy práctica en la clase es: cuando se obtiene un patch ``parcialmente correcto pero que no pasa'', ¿conviene seguir corrigiéndolo en la trayectoria actual o descartarlo y volver a muestrear? En realidad esto es \textbf{asignación de presupuesto de cómputo en tiempo de prueba}. Para cambios pequeños, volver a muestrear puede ser más barato; para modificaciones complejas de varios archivos, conservar el contexto y seguir corrigiendo puede ser mejor.

\begin{warningbox}{Error común: mirar solo la calidad de una trayectoria}
Si solo se analiza la mejor trayectoria, es fácil sobrestimar el sistema. En un despliegue real conviene observar al mismo tiempo:
\begin{itemize}
    \item el número promedio de trayectorias que consume cada éxito;
    \item si las trayectorias fallidas son recuperables;
    \item la proporción de fallos por tiempo agotado o por bucles.
\end{itemize}
\end{warningbox}

\subsection{La herramienta no es una función accesoria, sino el límite de la capacidad}
Como se plantea en la clase, la llamada a herramientas no es ``darle al modelo un poco de plugin'', sino que determina directamente el espacio de problemas resolubles. Sin una herramienta de exploración de código, el modelo casi no puede ubicar de manera estable los problemas de un repositorio grande; sin una herramienta de retroalimentación de ejecución, el modelo se queda por mucho tiempo en conjeturas semánticas; sin capacidad de depuración, en escenarios de seguridad la cadena de activación de vulnerabilidades es difícil de cerrar.

\begin{importantbox}{El significado sistémico de la llamada a herramientas}
El techo de capacidad de un Agente no lo determina solo el base model, sino que también lo define en conjunto ``el conjunto de herramientas $\times$ el formato de retroalimentación $\times$ la estrategia de control''. Quien imparte la clase llama a estos tres elementos un diseño a nivel de sistema que debe ajustarse en conjunto.
\end{importantbox}

\subsection{Resumen de la sección}
La definición de Agente en esta clase es clara: es un sistema de acción que puede iterar de manera sostenida, no un generador de respuestas largas. El cómputo dinámico, la retroalimentación de herramientas y las trayectorias recuperables constituyen la base de toda la discusión de ingeniería que sigue.

\section{Evolución del sistema de evaluación: de MBPP/HumanEval a SWE-Bench}

\subsection{El valor y la caducidad de las evaluaciones tempranas de código}
MBPP y HumanEval impulsaron en su momento la iteración rápida de los modelos de código, porque cumplían entonces tres condiciones clave: dificultad adecuada, calificación automática y una filtración aún no generalizada. Sutton señala directamente que hoy este tipo de tareas se parece más al ``MNIST de la evaluación de código'': todavía sirve para comparaciones básicas, pero difícilmente representa la capacidad real de ingeniería de software.

\begin{knowledgebox}{Interpretación de Pass@K y sus implicaciones de despliegue}
Pass@K mide la probabilidad de que ``al muestrear K veces, al menos una pase las pruebas''. Tiene valor para evaluar la diversidad, pero no siempre equivale a la experiencia del producto:
\begin{itemize}
    \item la finalización única en el IDE se acerca más a Pass@1;
    \item un agente con una etapa de verificación puede acercarse más a Pass@K;
    \item por lo tanto, si la métrica es válida depende de cómo se use realmente el sistema.
\end{itemize}
\end{knowledgebox}

\begin{warningbox}{Un desajuste de métricas provoca un rumbo equivocado}
Si el escenario del producto solo permite una salida, pero el equipo optimiza principalmente Pass@K, puede terminar entrenando el sistema para que ``el conjunto de muestras contenga la respuesta correcta'', en lugar de que ``la primera respuesta sea utilizable''. Este sesgo se vuelve muy evidente en la etapa de lanzamiento.
\end{warningbox}

\subsection{El avance de SWE-Bench}
SWE-Bench hizo avanzar la tarea de ``escribir funciones cortas'' a ``corregir issues reales en un repositorio de código real'': la entrada es la descripción del issue y el contexto del repositorio, la salida es un patch, y la verificación depende del conjunto de pruebas ya existente. En comparación con las evaluaciones tempranas, elevó de forma notable el realismo y obligó a los investigadores a enfrentar problemas de sistema como la navegación de código, la recuperación de contexto y la edición de múltiples archivos.

\begin{figure}[H]
\centering
\includegraphics[width=0.88\textwidth]{frame-02-swebench.jpg}
\caption{Captura de la clase donde se explica la forma de la tarea de SWE-Bench y su impacto \protect\footnotemark}
\end{figure}
\footnotetext{Intervalo de tiempo en el video: 00:14:55--00:15:40.}

\begin{importantbox}{Lo que SWE-Bench realmente cambió}
Lo que cambió no fue solo la puntuación, sino el paradigma de investigación y desarrollo. Los equipos comenzaron a invertir de manera sistemática en:
\begin{itemize}
    \item recuperación y navegación a nivel de repositorio de código;
    \item corrección impulsada por retroalimentación de ejecución;
    \item gestión de trayectorias de múltiples turnos;
    \item un ciclo cerrado de generación y verificación de patches.
\end{itemize}
\end{importantbox}

\subsection{Las limitaciones de SWE-Bench son igual de importantes}
Sutton lo dice con franqueza: el valor y las limitaciones de SWE-Bench coexisten. Las fuentes de datos se concentran en un número reducido de proyectos y ecosistemas de lenguajes específicos; el estilo de las descripciones de los issues no coincide del todo con la interacción real entre humano y máquina dentro del IDE; los commits históricos tienen el ruido de ``un commit que mezcla varios cambios''; el conjunto de pruebas no es un verifier perfecto, y pueden aparecer parches que ``sobreajustan las pruebas pero son semánticamente incorrectos''.

\begin{table}[H]
\centering
\begin{tabular}{p{3.4cm}p{5.0cm}p{5.0cm}}
\toprule
\textbf{Dimensión} & \textbf{Contribución de SWE-Bench} & \textbf{Sesgo potencial} \\
\midrule
Realismo de la tarea & Proviene de la cadena issue/patch de proyectos reales & Distribución de proyectos limitada, cobertura de dominios incompleta \\
Forma de verificación & Puede ejecutar pruebas automáticamente y formar una evaluación de ciclo cerrado & Las pruebas no representan completamente la corrección semántica \\
Entrada de lenguaje & Conserva el estilo de redacción real de los issues y las pistas de contexto & Existe una diferencia de distribución respecto a las instrucciones breves dentro del IDE \\
Valor de entrenamiento & Puede impulsar el diseño de retrieval, editing y execution & Los commits históricos pueden entrelazar múltiples objetivos de cambio \\
\bottomrule
\end{tabular}
\caption{El ``fuerte realismo'' y el ``fuerte sesgo'' de SWE-Bench coexisten}
\end{table}

\subsection{La vida útil de las evaluaciones y su renovación continua}
La clase plantea un juicio muy importante: \textbf{toda evaluación tiene una shelf life}. Al principio es difícil y no se ha filtrado, pero a los pocos años suele volverse fácil de resolver y de memorizar. El sistema de evaluación debe actualizarse de manera continua, o el sistema terminará optimizándose eficientemente sobre un objetivo caduco.

\begin{importantbox}{Una frase metodológica que se puede poner en práctica}
\textit{``If your numbers aren't good, you're hill climbing on the wrong thing.''}  
Traducido a una acción de ingeniería, esto significa: primero verificar si el objetivo de evaluación es real y si todavía tiene poder de discriminación, y solo después ajustar los parámetros del modelo o del harness.
\end{importantbox}

\subsection{Resumen de la sección}
De MBPP/HumanEval a SWE-Bench, la industria no se limitó a ``cambiar de tabla de clasificación'', sino que llevó la inteligencia de código del nivel de función al nivel de repositorio. Al mismo tiempo, la evaluación misma debe auditarse de manera continua, o una puntuación alta no equivaldrá a una alta utilidad práctica.

\section{Ingeniería del harness: codiseño de herramientas, interfaz y flujo de control}

\subsection{Por qué el harness es un ciudadano de primera clase}
Sutton dice a los 00:26 una frase que se cita con frecuencia en esta clase: \textit{``The design of the agent harness and the tools, you have to co-design it with whatever the model can do.''}
Esta frase significa que no debemos tratar al modelo como una caja negra a la que se le conectan herramientas por fuera, sino tratar la interfaz de las herramientas, el formato del prompt y la estructura del flujo de control como objetos de diseño del mismo nivel que el modelo.

\begin{importantbox}{Tres principios de ingeniería del Harness Engineering}
\begin{enumerate}
    \item La semántica de las acciones debe acercarse a la distribución que el modelo conoce, para reducir la ambigüedad en las llamadas;
    \item La granularidad de las herramientas debe ser lo bastante compacta para reducir los casos de ``llamar a cinco herramientas para hacer una sola cosa'';
    \item El formato de salida debe ser comprimible y rastreable, para evitar que el ruido sature el contexto.
\end{enumerate}
\end{importantbox}

\subsection{La perspectiva de Agent Computer Interface (ACI)}
En la clase se cita el concepto de ``ACI'' (Agent Computer Interface), con la intención de establecer una analogía con HCI: el diseño de interfaces persona-computadora tiene paradigmas establecidos, y la interfaz entre el agente y su entorno de ejecución necesita, de igual manera, un método sistemático. Por ejemplo, acciones como ``ir a la definición'' o ``buscar referencias por símbolo'' suelen ajustarse mejor a la ruta de pensamiento del modelo que un grep genérico, y se acercan más al flujo de trabajo de un IDE.

\begin{knowledgebox}{El principio de la ``ruta familiar'' en el diseño de herramientas}
El presentador retoma una idea parecida al ``Little Red Riding Hood principle'': conviene llevar al modelo, en la medida de lo posible, hacia los numerosos patrones de interacción que ya vio durante el entrenamiento posterior. Cuanto más se acerquen el nombre de la herramienta, sus parámetros y el texto que devuelve a los contextos de desarrollo habituales, más fácil será lograr llamadas exitosas de manera estable.
\end{knowledgebox}

\begin{figure}[H]
\centering
\includegraphics[width=0.88\textwidth]{frame-03-harness.jpg}
\caption{Imagen clave de la clase que enfatiza que ``el harness y las herramientas deben diseñarse en conjunto'' \protect\footnotemark}
\end{figure}
\footnotetext{Intervalo de tiempo en el video: 00:26:24--00:26:55.}

\subsection{El espectro del diseño del flujo de control: de la libertad a la restricción}
La clase ubica el flujo de control de los agentes de programación en un espectro continuo: en un extremo está ``dar shell más IDE, casi todo dinámico''; en el medio están los ciclos ReAct como el de SWE-agent; más hacia la derecha están los ensambles por etapas (primero recuperación de información, luego generación del parche); y en el extremo derecho están las máquinas de estados (state machine) explícitas, que restringen las acciones ejecutables en cada paso.

\begin{table}[H]
\centering
\begin{tabular}{p{3.2cm}p{4.4cm}p{3.8cm}p{2.6cm}}
\toprule
\textbf{Paradigma de flujo de control} & \textbf{Ventajas} & \textbf{Riesgos} & \textbf{Tareas idóneas} \\
\midrule
ReAct totalmente dinámico & Exploración flexible, se adapta a situaciones desconocidas & Deriva de la trayectoria, costo inestable & Localización y depuración abiertas \\
Ciclo de dos etapas ensambladas & Desacopla la recolección de información de la generación del parche & Un error en el límite entre etapas provoca fallas en cascada & Tareas de reparación de complejidad media \\
Ciclo restringido por máquina de estados & Alta controlabilidad, facilita la auditoría y la reproducción & Espacio de búsqueda limitado, puede pasar por alto soluciones & Escenarios de alto riesgo, procesos de cumplimiento \\
\bottomrule
\end{tabular}
\caption{El diseño del flujo de control es un equilibrio entre el riesgo de la tarea y la capacidad de exploración}
\end{table}

\begin{warningbox}{``Mientras más control, mejor'' es un malentendido}
Una restricción demasiado fuerte le quita al Agent la capacidad de recuperarse en rutas anómalas; una restricción demasiado débil produce alucinaciones en trayectorias largas. Lo correcto es configurar el flujo de control por capas, según la complejidad y el nivel de riesgo de la tarea, en lugar de buscar un solo paradigma.
\end{warningbox}

\subsection{Análisis de trayectorias: hay que revisar las muestras fallidas}
Antes de abordar el tema de los agentes de seguridad, Sutton repite un viejo principio del aprendizaje automático: \textit{``Always look at the data.''} Aquí ``data'' no se refiere al conjunto de entrenamiento, sino a las trayectorias de ejecución del Agent. Solo revisando una por una las rutas fallidas se puede saber si el problema fue un desajuste en la recuperación, un mal uso de una herramienta o un verificador que indujo a error.

\begin{importantbox}{Recomendación práctica: iteración del harness guiada por trayectorias}
\begin{itemize}
    \item En cada iteración de versión, tomar una muestra fija de trayectorias fallidas para revisarlas de manera manual;
    \item Traducir los patrones de falla más frecuentes en nuevas herramientas o nuevas restricciones;
    \item Consolidar la experiencia de ``ya lo arreglé una vez'' en la system instruction y en las plantillas de política.
\end{itemize}
\end{importantbox}

\begin{figure}[H]
\centering
\includegraphics[width=0.88\textwidth]{frame-04-control-loop.jpg}
\caption{Explicación centrada en el ciclo de flujo de control y de política, que enfatiza que el ciclo dinámico y la restricción programática no son opuestos \protect\footnotemark}
\end{figure}
\footnotetext{Intervalo de tiempo en el video: 00:38:10--00:38:55.}

\subsection{Resumen de la sección}
El núcleo de Harness Engineering es la ``coordinación del sistema'', no el ``ajuste fino del prompt''. La semántica de las herramientas, la estructura de la retroalimentación, las restricciones del flujo de control y la capacidad del modelo deben optimizarse en conjunto para poder mejorar de manera estable en tareas complejas.
\section{De la corrección de código a las tareas de seguridad: transferencia de capacidades y reconfiguración de escenarios}

\subsection{Diferencias estructurales entre las tareas de seguridad y las tareas de corrección comunes}
Las tareas de seguridad no se limitan a «hacer que las pruebas pasen»; también exigen responder tres preguntas más difíciles: si la vulnerabilidad es realmente activable, cuál es la condición que la activa y si puede reproducirse en distintas rutas. En comparación con la corrección de código en general, la búsqueda de vulnerabilidades depende más de la observación del comportamiento dinámico, del razonamiento sobre el flujo de datos entre funciones y de la verificación contrafactual.

\begin{knowledgebox}{El ciclo mínimo de tareas de un Agent de seguridad}
\begin{enumerate}
    \item Suponer que cierta ruta podría tener una condición de desbordamiento, acceso fuera de límites o falta de verificación;
    \item generar una entrada ejecutable o un script que la active;
    \item ejecutarla bajo un sanitizer o depurador y recopilar evidencia;
    \item si falla, corregir la suposición y seguir explorando hasta reproducir o refutar el caso.
\end{enumerate}
\end{knowledgebox}

\subsection{CTF/Enigma como escalón de transición}
La clase presenta las tareas de tipo CTF y Enigma como una capa intermedia: enfatizan más las cadenas de explotación y activación que los parches cotidianos dentro del IDE, pero siguen siendo más controlables que las vulnerabilidades de sistemas reales a gran escala. Este escalón tiene valor porque permite al equipo iterar sus herramientas y su flujo de control en un entorno evaluable.

\begin{importantbox}{Por qué empezar con CTF antes de pasar a vulnerabilidades reales}
El CTF permite verificar rápidamente, en un entorno reproducible, si «la herramienta sirve»; las vulnerabilidades reales, en cambio, están más influidas por el ruido del contexto, el entorno de dependencias y la complejidad del espacio de entradas. Pulir primero el ciclo básico del Agent y luego migrar hacia sistemas reales produce una tasa de éxito más alta.
\end{importantbox}

\subsection{Las herramientas de análisis dinámico entran al flujo principal}
En la segunda mitad del minuto 00:50, el presentador habló específicamente del papel del sanitizer como analizador dinámico. A diferencia de un juicio puramente estático, el sanitizer ofrece directamente, en tiempo de ejecución, la señal de «si se activó un comportamiento de memoria ilegal», lo que lo hace naturalmente apto para ser el oracle del Agent.

\begin{warningbox}{Depender solo de un clasificador estático lleva a caer en «parece una vulnerabilidad»}
Muchas funciones parecen peligrosas de forma local, pero la cadena de llamadas global puede ya incluir verificaciones de límites; a la inversa, algo que parece seguro localmente puede seguir siendo explotable al combinarse entre módulos. Sin verificación dinámica, tanto los falsos positivos como los falsos negativos aumentan de forma considerable.
\end{warningbox}

\subsection{Resumen de la sección}
Las tareas de seguridad no son una simple ampliación de las tareas de corrección de código, sino un cambio de la función objetivo. Buscar solo que las pruebas pasen no puede responder si «realmente existe una vulnerabilidad explotable»; es necesario introducir evidencia de ejecución dinámica y un ciclo de verificación más sólido.

\section{Por qué la detección de vulnerabilidades es difícil: contexto semántico, ruido en la evaluación y verificación dinámica}

\subsection{Caso del kernel: lo sospechoso a nivel local no implica que sea válido a nivel global}
El curso usó un ejemplo de código del kernel para explicar la dificultad de determinar vulnerabilidades: una función lee una longitud del encabezado de un paquete, asigna un búfer y copia datos; si la longitud no coincide con el payload real, puede producirse una escritura fuera de límites. La dificultad no está en el reconocimiento de patrones en sí, sino en \textbf{confirmar la ruta alcanzable, las restricciones del llamador y el flujo de datos entre funciones}.

\begin{figure}[H]
\centering
\includegraphics[width=0.88\textwidth]{frame-05-kernel-vuln.jpg}
\caption{Fragmento de la explicación de la vulnerabilidad del kernel: el riesgo de la misma función depende del contexto global de la llamada \protect\footnotemark}
\end{figure}
\footnotetext{Intervalo de tiempo en el video: 01:06:30--01:07:10.}

\begin{importantbox}{Juicio clave del ponente}
\textit{``This function on its own is not vulnerable or non-vulnerable.''}  
Clasificar una función de forma aislada suele no ser un problema bien definido. Que una vulnerabilidad se considere válida depende de si el llamador ya restringe la entrada, de si la ruta es alcanzable y de si el entorno de ejecución cumple las condiciones para desencadenarla.
\end{importantbox}

\subsection{La dificultad natural de evaluar con bases de datos de vulnerabilidades}
Las bases de datos públicas de vulnerabilidades (como el ecosistema CVE) parecen una fuente de datos ideal: tienen número, descripción y parche. Pero convertirlas directamente en una ``tarea de clasificación a nivel de función'' genera problemas serios: hasta dónde cortar el contexto de entrada, cómo definir los ejemplos positivos y negativos, si el parche corresponde a una sola causa raíz, y si la etiqueta es estable entre versiones.

\begin{warningbox}{``Descargar una base de datos de vulnerabilidades para entrenar un detector'' suele ser demasiado optimista}
La clase señaló explícitamente que esta es la ruta más natural, pero también la más fácil de errar. Si no se define primero un límite riguroso de la tarea, el modelo puede aprender el estilo del proyecto, los hábitos de los commits o frases de plantilla, en lugar de la semántica real de la explotabilidad.
\end{warningbox}

\subsection{La necesidad de una evaluación dinámica}
Sutton mencionó las competencias relacionadas con DARPA y la dirección de evaluación de seguridad de Meta; el punto en común es pasar de etiquetas estáticas a verificación ejecutable. Para la tarea de descubrimiento de vulnerabilidades, la evaluación dinámica se acerca más al objetivo real: activación por entrada, evidencia de ejecución, ubicación del fallo, reproducibilidad.

\begin{knowledgebox}{La relación complementaria entre fuzzing y los agentes}
\begin{itemize}
    \item El fuzzing es bueno para la exploración aleatoria de alto rendimiento;
    \item El agente es bueno para guiar por semántica y plantear hipótesis de ruta;
    \item Combinar ambos puede formar una búsqueda híbrida de ``cobertura amplia + ruta profunda''.
\end{itemize}
Esta es también la dirección que se enfatiza repetidamente en la última parte de la clase: la IA no sustituye al fuzzing, sino que empuja la búsqueda hacia las ramas lógicas profundas a las que el fuzzer difícilmente llega.
\end{knowledgebox}

\subsection{De la precisión estática a la tasa de éxito de la tarea}
En escenarios de seguridad, el indicador final no debería ser ``qué tanto se parece este código a una vulnerabilidad'', sino ``si se puede construir una entrada que la active y aportar evidencia auditable''. Esto traslada el foco de investigación de la clasificación de texto a \textbf{la tasa de éxito de tareas ejecutables}, algo que coincide de forma natural con el paradigma de agentes.

\begin{importantbox}{Implicación directa para el diseño de investigación}
Si en el futuro los artículos sobre agentes de seguridad solo reportan puntuaciones de clasificación estática, será difícil sustentar su valor de ingeniería; como mínimo deberían añadir indicadores como la tasa de activación dinámica, la estabilidad de reproducción y el costo de tratar los falsos positivos.
\end{importantbox}

\subsection{Resumen de la sección}
La dificultad de la detección de vulnerabilidades está en ``la dependencia del contexto más la verificación ejecutable'', no en ``la falta de reconocimiento de patrones''. Esto también explica por qué la seguridad necesita más sistemas basados en agentes, y no solo modelos de texto de código más potentes.
\section{Repaso del caso Big Sleep: el ciclo completo, de la hipótesis a la evidencia de fallo}

\subsection{Objetivo del proyecto y definición de la tarea}
Big Sleep (en el curso también se relaciona con Project Naptime) se centra en un objetivo que puede verificarse de forma explícita: dado un repositorio de código, generar una entrada que provoque un fallo detectado por el sanitizer. Este objetivo convierte «descubrir una vulnerabilidad» en un problema de búsqueda ejecutable, lo que permite tanto la evaluación automática como conservar la fidelidad a la práctica de ingeniería real.

\begin{figure}[H]
\centering
\includegraphics[width=0.88\textwidth]{frame-06-big-sleep.jpg}
\caption{Pantalla de presentación del proyecto Big Sleep: convertir el descubrimiento de vulnerabilidades en una tarea de disparo ejecutable \protect\footnotemark}
\end{figure}
\footnotetext{Intervalo de tiempo en el video: 01:14:58--01:15:35.}

\begin{importantbox}{El modelado clave de Big Sleep}
El objetivo no es que el modelo explique de palabra que ``aquí hay una vulnerabilidad'', sino que el sistema produzca una entrada reproducible que haga que el programa genere una anomalía objetiva bajo el sanitizer o el depurador. Dicho de otro modo, la conclusión debe ser reconocida por el entorno de ejecución.
\end{importantbox}

\subsection{Tres categorías centrales de herramientas}
El curso mostró tres categorías de herramientas de este sistema:
\begin{itemize}
    \item Herramienta de exploración de código: búsqueda por símbolo, salto a la definición, seguimiento de referencias;
    \item Herramienta de generación de entradas: que el Agente escriba un programa en Python para construir entradas complejas, en lugar de escribir a mano cadenas de bytes largas;
    \item Herramienta de ejecución con depuración: ejecutar el programa objetivo y devolver los valores de las expresiones correspondientes a watchpoints/breakpoints.
\end{itemize}

\begin{knowledgebox}{Por qué hacer que el Agente genere el programa de entrada en Python}
En escenarios de seguridad, las entradas suelen contener estructuras binarias, formatos anidados o campos de verificación. Que el Agente escriba un ``programa generador'' es más robusto que producir directamente bytes crudos, y también facilita las modificaciones iterativas posteriores.
\end{knowledgebox}

\begin{warningbox}{Tratar el depurador como una ventana de chat interactiva no es eficiente}
El diseño de herramientas mostrado en clase no se detiene en cada paso al estilo de la depuración manual paso a paso, sino que sigue el esquema de ``ejecutar el programa completo y devolver los valores de los puntos de observación alcanzados''. Esto reduce la sobrecarga de interacción y a la vez conserva la evidencia clave, lo que resulta más adecuado para un ciclo automatizado.
\end{warningbox}

\subsection{Resolución de problemas mediante trayectorias}
La trayectoria que muestra el curso es muy representativa: primero se explora el código y se formula una hipótesis de vulnerabilidad, después se produce un informe de la vulnerabilidad junto con una idea de explotación; a continuación se genera la entrada y se ejecuta; si no se produce el fallo, se atribuye la causa del fracaso, se corrige la entrada y se vuelve a ejecutar, hasta provocar la anomalía. Este proceso refleja la ventaja central del Agente: \textbf{el fracaso no es el final, sino la señal de supervisión para la siguiente ronda de búsqueda}.

\begin{figure}[H]
\centering
\includegraphics[width=0.88\textwidth]{frame-08-wrapup.jpg}
\caption{Síntesis de cierre de la clase sobre el estado actual y las perspectivas de los Agentes de seguridad \protect\footnotemark}
\end{figure}
\footnotetext{Intervalo de tiempo en el video: 01:25:45--01:26:20.}

\begin{importantbox}{Un hecho clave a nivel de resultados}
El curso reveló un resultado del mundo real: se descubrió y se divulgó una vulnerabilidad real en SQLite. Más importante aún, el proyecto también mostró capacidad de variant analysis, es decir, buscar variantes similares dentro del mismo proyecto a partir de patrones de vulnerabilidades ya conocidos.
\end{importantbox}

\subsection{Comparación con el fuzzing tradicional}
El expositor mencionó un detalle muy convincente: incluso conociendo de antemano la ubicación de la vulnerabilidad, al intentar encontrarla en sentido inverso con un fuzzer genérico, una ejecución prolongada (el curso menciona un orden de magnitud de unas 150 horas) sigue teniendo dificultades para dar directamente con la misma vulnerabilidad. Esta comparación sugiere que la guía semántica del Agente puede aportar una ganancia en la exploración de rutas profundas.

\begin{knowledgebox}{El valor potencial de la búsqueda agentic}
\begin{itemize}
    \item Puede aprovechar conocimiento previo (vulnerabilidades históricas, semántica de los parches) para realizar una búsqueda dirigida;
    \item Puede reescribir la estrategia de generación de entradas después de un fracaso;
    \item Puede combinarse con la retroalimentación de la depuración para reducir el espacio de rutas sospechosas.
\end{itemize}
\end{knowledgebox}

\begin{warningbox}{No hay que extrapolar el éxito de un solo caso a una ``superioridad total''}
Big Sleep demostró la viabilidad, pero eso no significa que todos los proyectos y todos los tipos de vulnerabilidades ya puedan automatizarse. La conclusión de la clase es que ``apenas está empezando y hay mucho margen'', no que ``el problema ya esté resuelto''.
\end{warningbox}

\subsection{Resumen de la sección}
El valor de Big Sleep está en que ejemplifica un ciclo completo: el objetivo es ejecutable, las herramientas son invocables, la trayectoria es iterable y el resultado es verificable. Hace avanzar a los Agentes de seguridad, de la prueba de concepto a una práctica de ingeniería reproducible.
\section{Hoja de ruta de implementación en ingeniería: de los paradigmas de la clase al experimento reproducible}

\subsection{Primero se define el «objeto optimizable», después se optimiza el modelo}
Lo más útil de trasladar esta clase al desarrollo cotidiano es lo siguiente: \textbf{primero hay que definir la tarea como un objeto ejecutable, verificable e iterable, y solo después hablar de optimizar el modelo}. Muchos equipos fracasan no porque el modelo no sea lo bastante potente, sino porque la definición de la tarea está incompleta; por ejemplo, solo cuentan con un objetivo en lenguaje natural sin un verifier explícito, o solo con la puntuación final sin diagnóstico a nivel de trayectoria.

\begin{importantbox}{La unidad experimental mínima de un Code Agent de seguridad}
Un experimento reproducible debe incluir al menos seis componentes:
\begin{enumerate}
    \item especificación de la tarea: entrada, salida, restricciones, tiempo límite;
    \item entorno de ejecución: contenedor reconstruible y versiones de las dependencias;
    \item interfaz de herramientas: exploración de código, ejecución, depuración, generación de entradas;
    \item verificador: conjunto de pruebas o sanitizer/oracle;
    \item registro de trayectoria: thought/tool/observation en cada ronda;
    \item script de evaluación: ejecución por lotes y agregación estadística.
\end{enumerate}
\end{importantbox}

Si faltan los puntos 5 y 6, el sistema suele quedar en la situación de ``ver solo el resultado, sin ver el proceso'', lo que dificulta reproducir los problemas. Aquí es donde realmente aterriza lo que Sutton dice con ``always look at the data'': el equipo debe poder tratar las trayectorias fallidas como un activo de datos de primer orden y analizarlas de manera continua.

\begin{knowledgebox}{Estrategia de puente entre la clase y la ingeniería}
La estructura en tres partes del curso puede trasladarse a un flujo de desarrollo de tres capas:
\begin{itemize}
    \item capa base: evaluación de tareas de código (MBPP/HumanEval/SWE-Bench) para establecer las capacidades fundamentales;
    \item capa de sistema: harness, herramientas y flujo de control para mejorar la estabilidad en tareas complejas;
    \item capa de aplicación: verificación dinámica y cadena de evidencia en tareas de seguridad para comprobar el valor real.
\end{itemize}
\end{knowledgebox}

\subsection{Ritmo de iteración: descomponer los cuellos de botella del sistema por tipo de fallo}
En clase se mencionó varias veces la ``optimización combinatoria'', que en ingeniería puede llevarse un paso más allá mediante la clasificación de fallos. En cuanto las causas de los fallos se clasifican de forma estable, las prioridades de ajuste quedan mucho más claras, en lugar de reemplazar el modelo a ciegas o apilar prompts.

\begin{table}[H]
\centering
\begin{tabular}{p{3.2cm}p{4.6cm}p{4.6cm}}
\toprule
\textbf{Tipo de fallo} & \textbf{Síntoma típico} & \textbf{Dirección prioritaria de corrección} \\
\midrule
Fallo de recuperación & Explora repetidamente archivos irrelevantes, la localización de la ruta se desvía & Mejorar la búsqueda a nivel de símbolo, incorporar navegación por referencias cruzadas \\
Fallo de ejecución & Formato incorrecto en la invocación de comandos, dependencias de entorno faltantes & Restringir los parámetros de las herramientas, fijar una línea base del contenedor \\
Fallo de verificación & Las pruebas pasan pero hay error semántico, o se juzga mal la causa del fallo & Reforzar el verifier, añadir repetición de contraejemplos \\
Fallo de estrategia & En trayectorias largas se autorrefuerzan hipótesis erróneas & Reiniciar el flujo de control, ciclos por etapas, restricciones de máquina de estados \\
Fallo de costo & La tasa de éxito mejora, pero el consumo de tokens o la latencia son inaceptables & Gestión del presupuesto, estrategias de parada temprana, enrutamiento de modelos por niveles \\
\bottomrule
\end{tabular}
\caption{Mapeo de los modos de fallo del Agente a direcciones de mejora del sistema que se puedan poner en práctica}
\end{table}

\begin{warningbox}{Revisar solo «casos de éxito» induce al equipo a error}
Si en las reuniones de revisión solo se muestra la best trajectory, el sistema terminará optimizándose para ser ``ocasionalmente impresionante''. Para mejorar la fiabilidad en producción es necesario revisar de manera sistemática las muestras de fallo de cola larga, en especial las trayectorias que ``parecen a punto de tener éxito pero terminan colapsando''.
\end{warningbox}

\subsection{Gobernanza del flujo de control: el equilibrio entre eficiencia de exploración y capacidad de auditoría}
Según el espectro continuo que presenta el curso, el flujo de control no consiste en tomar partido por un extremo, sino en configurarlo según el escenario. La fase de exploración puede inclinarse hacia lo dinámico, y la fase de convergencia hacia lo restringido. En las tareas de seguridad, además, se deben cumplir requisitos de auditoría: que cada acción clave se pueda rastrear y que cada conclusión se pueda reproducir.

\begin{importantbox}{Flujo de control de doble modo recomendado}
\begin{itemize}
    \item modo de exploración: muestreo alto, múltiples rutas, permisos amplios de herramientas, para descubrir posibles rutas de vulnerabilidad;
    \item modo de verificación: muestreo bajo, restricciones estrictas, puntos de observación fijos, para reproducir y consolidar evidencia.
\end{itemize}
Ambos modos comparten el mismo objetivo de la tarea, pero difieren en la estrategia de recursos y en los requisitos de interpretabilidad.
\end{importantbox}

\begin{knowledgebox}{Reglas de gestión de presupuesto que se pueden poner en práctica}
El presupuesto en tiempo de prueba puede dividirse en tres capas:
\begin{itemize}
    \item presupuesto por trayectoria: limita el número máximo de pasos y de ejecuciones de cada trayectoria;
    \item presupuesto por tarea: limita el número total de trayectorias y la duración total de cada ejemplo;
    \item presupuesto por lote: limita el total de horas de GPU/CPU de toda la ronda de evaluación.
\end{itemize}
Cuando alguna capa alcanza su umbral, el sistema debe generar una razón estructurada de ``por qué se detuvo'', para evitar que el tiempo de espera agotado sea una caja negra.
\end{knowledgebox}

\subsection{Restricciones adicionales para la implementación en seguridad: divulgación responsable y riesgo de doble uso}
El caso Big Sleep demuestra que ``descubrir una vulnerabilidad'' ya es viable en la práctica, por lo que, al ponerlo en producción, no basta con hablar de métricas técnicas. Un Agente de seguridad necesita límites de responsabilidad integrados, entre ellos un proceso de gestión de vulnerabilidades, un ritmo de divulgación pública y salvaguardas que impidan que la capacidad se reconvierta directamente en automatización de ataques.

\begin{warningbox}{El riesgo de doble uso debe gobernarse desde el principio}
Cuanto más automatizado sea un sistema de descubrimiento de vulnerabilidades, más necesario es convertir en flujo predeterminado el control de acceso, el sandbox de ejecución, la anonimización de resultados y la aprobación de divulgación, en lugar de tratarlos como parches de última hora en el proyecto. De lo contrario, el éxito técnico se convierte directamente en riesgo de cumplimiento normativo.
\end{warningbox}

\begin{importantbox}{Lista mínima de gobernanza de seguridad para la práctica}
\begin{enumerate}
    \item incluir en lista blanca todos los repositorios objetivo y entornos de ejecución;
    \item ejecutar por defecto las entradas de tipo exploit-like solo en un entorno de ejecución aislado;
    \item al detectar una posible vulnerabilidad, verificarla internamente antes de iniciar el proceso de divulgación;
    \item distinguir en el informe entre ``evidencia'' e ``inferencia'', para evitar conclusiones exageradas;
    \item al publicar hacia el exterior, eliminar los detalles sensibles que permitan reproducir directamente la cadena de ataque.
\end{enumerate}
\end{importantbox}

\subsection{Oportunidades de investigación: preguntas que vale la pena abordar a continuación}
Al final, Sutton subraya que este es un ``wide open area''; si se traduce en una agenda de investigación, puede concentrarse en cuatro tipos de problemas: verifiers dinámicos más sólidos, generalización entre repositorios más estable, búsqueda de trayectorias largas con menor costo y generación automática de evidencia más confiable. En conjunto, estos determinan si un Agente de seguridad puede pasar de un éxito puntual a una producción estable y a escala.

\begin{table}[H]
\centering
\begin{tabular}{p{3.0cm}p{4.4cm}p{4.8cm}}
\toprule
\textbf{Dominio del problema} & \textbf{Cuello de botella actual} & \textbf{Posible dirección de avance} \\
\midrule
Verificador dinámico & La señal de activación es inestable y el costo de descartar falsos positivos es alto & Combinar sanitizer + aserciones semánticas + ejecución diferencial \\
Generalización entre proyectos & Las herramientas y los prompts dependen en gran medida de las convenciones de un solo repositorio & Protocolo ACI transferible e interfaz de normalización de tareas \\
Eficiencia de búsqueda & Las trayectorias largas tienen un costo alto y la retroalimentación de los fallos es escasa & Planificación por capas, repetición de experiencia, reutilización de la memoria de fallos \\
Automatización de evidencia & El informe carece de legibilidad y de capacidad de auditoría & Plantillas estructuradas de evidencia de vulnerabilidad y artifacts reproducibles \\
\bottomrule
\end{tabular}
\caption{Mapa de oportunidades de investigación a corto y mediano plazo para Agentes de seguridad}
\end{table}

\subsection{Resumen de la sección}
Trasladar el método de la clase a la práctica de ingeniería depende de dos cosas: primero, toda optimización debe girar en torno a un ciclo cerrado ejecutable; segundo, toda capacidad debe ir acompañada de mecanismos de auditoría y gobernanza. Solo cuando se cumplen a la vez ``poder hacerlo'' y ``poder controlarlo'', un Agente de seguridad puede entrar de manera sostenida en un flujo de producción real.

\section{Síntesis y ampliación}

\subsection{Repaso del marco general de la clase}
Esta clase enlazó tres líneas: la línea de evaluación (qué métricas vale la pena optimizar), la línea de sistemas (cómo diseñar el harness y las herramientas) y la línea de seguridad (por qué es necesario verificar de forma dinámica). Las tres líneas convergen finalmente en la misma conclusión: \textbf{la competitividad de un Agent proviene de la calidad del ciclo cerrado del sistema, no de una capacidad textual de un solo intento}.

\begin{table}[H]
\centering
\begin{tabular}{p{3.1cm}p{4.4cm}p{4.2cm}p{2.6cm}}
\toprule
\textbf{Tema} & \textbf{Conclusión de esta clase} & \textbf{Práctica aplicable} & \textbf{Desafío pendiente} \\
\midrule
Diseño de evaluación & La evaluación orienta el rumbo de la investigación y el desarrollo, pero tiene una vida útil & Actualizar continuamente la distribución de tareas, calibrar con escenarios reales & Resistir la fuga de datos y el sobreajuste \\
Arquitectura del Agent & ReAct más retroalimentación de herramientas es la unidad básica & Diseño del flujo de control por capas, las trayectorias fallidas deben auditarse & Equilibrio entre costo y estabilidad \\
Ingeniería del harness & La interfaz de herramientas debe diseñarse junto con la capacidad del modelo & Pensamiento ACI, acciones compactas, salidas comprimibles & Compromiso entre generalidad y especialización \\
Tareas de seguridad & Se debe pasar del juicio estático a la verificación dinámica & sanitizer/debugger entran al bucle principal & Generalización entre proyectos y control de falsos positivos \\
Práctica de Big Sleep & El descubrimiento de vulnerabilidades reales ya es alcanzable & Iterar usando la evidencia de activación como función objetivo & Extenderse a pilas de sistemas más complejas \\
\bottomrule
\end{tabular}
\caption{Tabla resumen de las conclusiones centrales de Lecture 08}
\end{table}

\subsection{Further Reading}
\begin{itemize}
    \item Las líneas base centrales mencionadas en la conferencia de Sutton: MBPP, HumanEval, SWE-Bench, SWE-agent, AutoCodeRover.
    \item Evaluación y datos de seguridad: bases de datos públicas de vulnerabilidades (el ecosistema CVE), benchmarks de CTF, la línea de evaluación de seguridad de Meta, competencias relacionadas con DARPA.
    \item Dirección de ingeniería: Agent Computer Interface, análisis de fallos a nivel de trayectoria, diseño de harness que prioriza la verificación dinámica.
    \item Práctica de investigación en seguridad: materiales de divulgación pública de Big Sleep / Project Naptime, así como el proceso de divulgación de vulnerabilidades en colaboración con Project Zero.
\end{itemize}

\subsection{Resumen de la sección}
El juicio de Sutton en la etapa final puede tomarse directamente como hoja de ruta: \textit{``This is a wide open area. We're just getting started.''} Para los investigadores y los equipos de ingeniería, el siguiente punto de énfasis no es hacer otra ronda de puntuación superficial en benchmarks, sino integrar al Agent en flujos de trabajo de seguridad más reales, más auditables y más reproducibles.

\end{document}
